\section{Computations and results}
\label{sec:oneloopcalc}
\label{sec:oneloop}
While the matching coefficients for the fields can be gauge-dependent, it is important to note that the combinations leading to the local current operators under consideration in this work are gauge invariant. Therefore, we have chosen to employ the Feynman gauge consistently throughout our calculations. Our computations are conducted in the context of Euclidean spacetime. In Euclidean spacetime, the renormalized Lagrangian $\mathcal{L}_{h_v}$ is given by
\begin{eqnarray}
\label{eq:extendLRenEuclidean}
\mathcal{L}_{h_v}^{\rm E} =Z_{h_v}\bar{h}_v (-iv\cdot \partial + i\delta m) h_v + gZ_gZ_A^{\frac{1}{2}}Z_{h_v} \bar{h}_v v\cdot A^{a}T^a h_v,
\end{eqnarray}
where the superscript ``E” indicates Euclidean spacetime. Consequently, 
the Feynman rules for the ``heavy” auxiliary field propagator, with momentum $k$ flowing along the ``heavy” auxiliary field line,  and the  interaction vertex are given by 
$\frac{1}{k\cdot v -i0}$ and $-gv^{\mu}T^a$, respectively. The $-i0$ prescription in Euclidean spacetime plays key role in obtaining correct linear divergences and ``mass correction”  $i\delta m$.

We decompose a vector into parallel and transverse components according to ($v^2=1$ in Euclidean spacetime in our settings)
\begin{eqnarray}
q_{\|}^{\mu} = (v\cdot q) v^{\mu}, \, \, \, \, q_{\perp}^{\mu} = q^{\mu} - v\cdot q v^{\mu}.
\end{eqnarray}
We use the following Schwinger parametrization ($\alpha$-representation)
\begin{eqnarray}
\label{eq:alphax}
\frac{1}{A^n} = \frac{1}{\Gamma(n)}\int_0^{+\infty} d\alpha\, \alpha^{n-1} e^{-A\alpha},
\end{eqnarray}
for the QCD propagator denominators. For the eikonal propagator, we use another version of Schwinger parametrization
\begin{eqnarray}
\label{eq:alphay}
\frac{1}{(k\cdot v- i0)^n} = \frac{(i)^n}{\Gamma(n)}\int_0^{+\infty} d\alpha\, \alpha^{n-1} e^{-i(k\cdot v)\alpha},
\end{eqnarray}
with $k\cdot v$  real. 


\subsection{$\delta m$ \& $\zeta_{h_v}$}
\begin{figure}
\centering
\includegraphics[width =0.32\textwidth]{images/Wilsonlineself.pdf}
\caption{One-loop self-energy Feynman diagram for the ``heavy” auxiliary field in gradient-flow formalism. The dashed arrow line denotes the ``heavy” auxiliary field, the filled squares represent the flowed $B_{\mu}$ field at flow time $t$.}
\label{fig:heavyself}
\end{figure}
The calculation of the self-energy of the ``heavy” auxiliary field at one-loop in gradient flow is similar to that of the heavy quark field in HQET.
Setting the external momentum $p$ flowing from left to right, the one-loop self-energy diagram shown in figure \ref{fig:heavyself} gives the amplitude
\begin{eqnarray}
\mathcal{M}_{(a)} = g^2\tilde{\mu}^{2\epsilon}C_F\int\frac{d^dk}{(2\pi)^d}\frac{1}{\left(k\cdot v + p\cdot v - i0\right) k^2 } e^{-2tk^2},
\end{eqnarray}
where $\tilde{\mu}^{2\epsilon}= \left(\frac{\mu^2 e^{\gamma_E}}{4\pi} \right)^{\epsilon}$ and we have chosen  the ``heavy” auxiliary field to be in fundamental representation. For the adjoint representation case, we just need to replace $C_F$ with $C_A$ in the above expression.

Summing up geometric series of the one-loop self-energy diagrams of the ``heavy” auxiliary field propagator gives $1/(p\cdot v - \mathcal{M}_{(a)})$, which indicates that the one-loop ``mass correction” is given by $(-\mathcal{M}_{(a)})$. For simplicity, we set $p=0$ in $\mathcal{M}_{(a)}$, yielding
\begin{eqnarray}
\label{eq:deltamstart}
-\mathcal{M}_{(a)}|_{p=0}= -g^2\tilde{\mu}^{2\epsilon}C_F\int\frac{d^dk}{(2\pi)^d}\frac{1}{\left(k\cdot v- i0\right)k^2} e^{-2tk^2}.
\end{eqnarray}
Note that the above expression is not odd under $k\cdot v\rightarrow -k\cdot v$ and the $k\cdot v$ integration does not give a vanishing result  because the sign of the $i0$ prescription matters. 

 To derive the matching coefficient of the  ``heavy” auxiliary field at one-loop, we differentiate $\mathcal{M}_{(a)} $ with respect to $p\cdot v$,  and set $p=0$ for simplicity, resulting in
\begin{eqnarray}
\label{eq:Zhvstart}
\frac{\partial\mathcal{M}_{(a)}}{\partial (p\cdot v)}|_{p=0}= - g^2\tilde{\mu}^{2\epsilon}C_F\int\frac{d^dk}{(2\pi)^d}\frac{1}{\left(k\cdot v- i0\right)^2k^2} e^{-2tk^2}.
\end{eqnarray}

Now, we explicitly evaluate eq.~(\ref{eq:deltamstart}) and eq.~(\ref{eq:Zhvstart}) using the Schwinger parameterization in eq.~(\ref{eq:alphax}) and eq.~(\ref{eq:alphay}). We find (we set $\epsilon=0$ whenever the integral is finite)
\begin{eqnarray}
\label{eq:deltamcalc0}
-\mathcal{M}_{(a)}|_{p=0} &=& -ig^2C_F\int\frac{d^4k}{(2\pi)^4}\int_0^{+\infty}dx\int_0^{+\infty}dy\, e^{-(2t+x)k^2- i(k\cdot v)y},\\
\label{eq:Zhvcalc0}
\frac{\partial\mathcal{M}_{(a)}}{\partial (p\cdot v)}|_{p=0} &=& g^2\tilde{\mu}^{2\epsilon}C_F\int\frac{d^dk}{(2\pi)^d}\int_0^{+\infty}dx\int_0^{+\infty}dy\, y\, e^{-(2t+x)k^2- i(k\cdot v)y},
\end{eqnarray}
Performing the $k$ momentum integration gives 
\begin{eqnarray}
\label{eq:deltamcalc1}
-\mathcal{M}_{(a)}|_{p=0} &=& -i\frac{\alpha_s}{4\pi}C_F\int_0^{+\infty}dx\int_0^{+\infty}dy\, \frac{1}{\left(2t +x\right)^2}e^{-\frac{y^2}{4\left(2t+x\right)}},\\
\label{eq:Zhvcal1}
\frac{\partial\mathcal{M}_{(a)}}{\partial (p\cdot v)}|_{p=0} &=& \frac{g^2\tilde{\mu}^{2\epsilon}C_F}{(4\pi)^{d/2}} \int_0^{+\infty}dx\int_0^{+\infty}dy\, y\, \frac{1}{\left(2t +x\right)^{d/2}}e^{-\frac{y^2}{4\left(2t+x\right)}}.
\end{eqnarray}
Integrating out $x, y$, we obtain the results 
\begin{eqnarray}
\label{eq:deltamresult0}
-\mathcal{M}_{(a)}|_{p=0}&=& -i\frac{\alpha_s}{4\pi}C_F\frac{\sqrt{2\pi}}{\sqrt{t}},\\
\label{eq:Zhvresult0}
\frac{\partial\mathcal{M}_{(a)}}{\partial (p\cdot v)}|_{p=0} &=& - 2\times\frac{\alpha_sC_F}{4\pi}\left[\frac{1}{\epsilon_{\rm IR}}+\log\left(2\mu^2 te^{\gamma_E}\right)\right].
\end{eqnarray} 

Repeating the same calculation using dimensional regularization with $t=0$ gives vanishing $\mathcal{M}_{(a),\, t=0}|_{p=0}$ and 
\begin{eqnarray}
\frac{\partial\mathcal{M}_{(a),\, t=0}}{\partial (p\cdot v)}|_{p=0}=
2\frac{\alpha_sC_F}{4\pi}\left[\frac{1}{\epsilon_{\rm UV}}-\frac{1}{\epsilon_{\rm IR}}\right].
\end{eqnarray}
This confirms that we just need to subtract the IR poles from gradient flow results for the matching from the gradient-flow scheme to the $\overline{\rm MS}$ scheme when external momenta are set to zero.

Therefore, we obtain the following results (notice that the ``mass correction” in the renormalized Lagrangian is denoted as $i\delta m$)
\begin{eqnarray}
\label{eq:deltamresult}
\delta m &=&  -\frac{\alpha_s}{4\pi}C_R\frac{\sqrt{2\pi}}{\sqrt{t}}+ O(\alpha_s^2),\\
\label{eq:Zhvresult}
\zeta_{h_v}  &=& 1- \frac{\alpha_sC_R}{4\pi}\log\left(2\mu^2 te^{\gamma_E}\right) + O(\alpha_s^2),
\end{eqnarray}
where the $C_R =C_F$ for fundamental representation case and $C_R=C_A$ for adjoint representation case. The logarithmic dependence on $t$ in the above result of $\zeta_{h_v}$ is consistent with the renormalization constant of $h_v$ calculated in \cite{Craigie:1980qs}.

In the following matching calculations, we drop the calculations of un-flowed diagrams since they can be easily obtained from the amplitudes of the  corresponding flowed diagrams by setting $t=0$ from the beginning, which lead to results proportional to $\frac{1}{\epsilon_{\rm UV}} - \frac{1}{\epsilon_{\rm IR}}$ and obviously match the IR poles from flowed diagrams. We have explicitly checked that  the coefficients of $-\log(2\mu^2te^{\gamma_E})$ in the matching coefficients calculated below and above align with those of the UV poles  from the corresponding un-flowed operators at one-loop level , which can be served as a check point of our calculation.

%%%%%%%%%%%%%%%%%%%%%%%%%%%%%%%%%%%%%
\subsection{$\mathring{\zeta}_{\psi}$}
\begin{figure}
\centering
\includegraphics[width =0.72\textwidth]{images/quarkself.pdf}
\caption{One-loop quark self-energy Feynman diagrams in the gradient-flow formalism (mirror diagrams are implicit). The solid line is the fermion line, the double solid line is the flowed fermion line, the open circle represents  the flow vertex. }
\label{fig:quarkself}
\end{figure}
The Feynman diagrams for the one-loop corrections of the quark self-energy in the gradient-flow formalism are displayed in figure \ref{fig:quarkself}. A similar calculation was performed in ref.~\cite{Rizik:2020naq}, here we repeat the calculation for completeness, adhering to our strategy of splitting the matching coefficients into several parts that can be calculated separately. 

Diagram $(c_1)$ is the same as the one in ordinary QCD in the small flow-time limit, which gives the contribution 
\begin{eqnarray}
\label{eq:c1}
\frac{1}{i}\frac{\partial\mathcal{M}_{(c_1)}}{\partial \slashed{p}}|_{\slashed{p}=0}&=& -\frac{\alpha_sC_F}{4\pi}\left[\frac{1}{\epsilon_{\rm UV}}-\frac{1}{\epsilon_{\rm IR}}\right].
\end{eqnarray}
For diagrams $(c_2)$ and  $(c_3)$, we have 
\begin{eqnarray}
\label{eq:c2}
\mathcal{M}_{(c_2)}&=& 2\times (i)^2g^2\tilde{\mu}^{2\epsilon}C_F\int_0^tds\, \int\frac{d^dk}{(2\pi)^d}\frac{(-2\slashed{k})\slashed{k}}{\left(k^2 \right)^2}e^{-sk^2}\nonumber\\
&=&2\times\frac{\alpha_sC_F}{4\pi}\left[\frac{1}{\epsilon_{\rm UV}}+1+\log\left(2\mu^2 te^{\gamma_E}\right)\right],
\end{eqnarray}
and
\begin{eqnarray}
\label{eq:c3}
\mathcal{M}_{(c_3)}&=&2\times \frac{1}{2}g^2\tilde{\mu}^{2\epsilon}(-2dC_F)\int_0^tds\, \int\frac{d^dk}{(2\pi)^d}\frac{1}{k^2}e^{-2sk^2}\nonumber\\
&=&-4\times\frac{\alpha_sC_F}{4\pi}\left[\frac{1}{\epsilon_{\rm UV}}+\frac{1}{2}+\log\left(2\mu^2 te^{\gamma_E}\right)\right].
\end{eqnarray}

Diagram $(c_4)$ and its corresponding mirror diagram contribute proportionally to the external fermion momentum (with the gradient flow gauge parameter $\kappa=1$) and subsequently vanish upon setting the external momentum to zero. Diagram $(c_5)$ yields a null contribution due to its inherent oddness under the gluon loop momentum transformation $k\rightarrow -k$.

Adding up $\frac{1}{i}\frac{\partial\mathcal{M}_{(c_1)}}{\partial \slashed{p}}|_{\slashed{p}=0}$, $\mathcal{M}_{(c_2)}$ and $\mathcal{M}_{(c_3)}$, removing the UV divergences by introducing the ringed fermion field,  and subtracting the IR divergences through matching procedure, we obtain the one-loop matching coefficient for the matching from the ringed fermion field $\mathring{\chi}$ to the $\overline{\rm MS}$ renormalized fermion field $\psi$,
\begin{eqnarray}
\label{eq:quarkselfresult}
\mathring{\zeta}_{\psi}= 1 +  \frac{1}{2} \times\frac{\alpha_s}{4\pi}C_F\left[\log\left(2\mu^2 te^{\gamma_E}\right)-\log(432)\right] + O(\alpha_s^2).
\end{eqnarray}
%%%%%%%%%%%%%%%%%%%%%%%%%%%%%%%%%%%%%%%%%%%%
\subsection{$\zeta_{\psi h_v}$}
\begin{figure}
\centering
\includegraphics[width =0.62\textwidth]{images/qvertex.pdf}
\caption{Feynman diagrams of  the one-loop corrections for the $\bar{\psi}h_v$ vertex in gradient-flow formalism. The cross circle represents the $\bar{\psi}h_v$ vertex. The fermion field connected to the cross circle is flowed at  flow time $t$.}
\label{fig:qvertex}
\end{figure}
The one-loop correction Feynman diagrams for the $\bar{\psi}h_v$ vertex in gradient-flow formalism are given in figure \ref{fig:qvertex}. 
Setting the quark external momentum to be zero, we have the amplitude for diagram $(b_1)$,
\begin{eqnarray}
\label{eq:ZJstart}
\mathcal{M}_{(b_1)}= -g^2\tilde{\mu}^{2\epsilon}C_F\int\frac{d^dk}{(2\pi)^d}\frac{\left(\gamma\cdot v\right) \slashed{k}}{\left(-k\cdot v- i0\right)\left(k^2\right)^2} e^{-2tk^2}.
\end{eqnarray}
Obviously, the transverse components of $\slashed{k}$ give vanishing loop integration, thus we have
\begin{eqnarray}
\label{eq:ZJresult0}
\mathcal{M}_{(b_1)}&=&g^2\tilde{\mu}^{2\epsilon}C_F\int\frac{d^dk}{(2\pi)^d}\frac{1}{\left(k^2\right)^2} e^{-2tk^2}\nonumber\\
&=&\frac{\alpha_sC_F}{4\pi}\left[-\frac{1}{\epsilon_{\rm IR}}-\log\left(2\mu^2 te^{\gamma_E}\right) -1\right].
\end{eqnarray}

Diagram $(b_2)$ gives contribution proportional to the external  fermion momentum (with the  gradient flow gauge parameter $\kappa=1$) and subsequently vanishes upon setting the external momentum to zero.
Therefore, we have the one-loop result of  $\zeta_{\psi h_v}$ after matching to the $\overline{\rm MS}$ scheme,
\begin{eqnarray}
\label{eq:Zqhvresults}
\zeta_{\psi h_v}= 1 -\frac{\alpha_sC_F}{4\pi}\left[\log\left(2\mu^2 te^{\gamma_E}\right)+1\right] + O(\alpha_s^2).
\end{eqnarray} 
%%%%%%%%%%%%%%%%%%%%%%%%%%%%%%%%%%%%%%%%%%%%%%%%
\subsection{$\zeta_A$}
\begin{figure}
\centering
\includegraphics[width =0.72\textwidth]{images/gluonself.pdf}
\caption{Feynman diagrams for the one-loop gluon self-energy in gradient-flow formalism (mirror diagrams are implicit). The  blob in diagram $(d_1)$ represents the gluon vacuum  polarization. The single curvy lines are gluon lines, while the double curvy lines are the flowed gluon lines.}
\label{fig:gluonself}
\end{figure}
The Feynman diagrams for the one-loop corrections of gluon self-energy in the gradient-flow formalism are given by figure \ref{fig:gluonself}. For simplicity, we set the external momentum to be zero while extracting the corresponding contribution. Diagram $(d_1)$ is just the conventional gluon vacuum polarization, which gives the contribution in the small flow-time limit,
\begin{eqnarray}
\label{eq:d1}
\mathcal{M}_{(d_1)}&=&\frac{\alpha_s}{4\pi}\delta^{\mu\nu}\left(\frac{5}{3}C_A - \frac{2}{3}n_f\right)\left[\frac{1}{\epsilon_{\rm UV}}-\frac{1}{\epsilon_{\rm IR}}\right].
\end{eqnarray}

Diagrams $(d_2)$ and  $(d_3)$ give the contributions
\begin{eqnarray}
\label{eq:d2}
\mathcal{M}_{(d_2)}
&=& 4g^2C_A\tilde{\mu}^{2\epsilon}\int_0^t ds \int\frac{d^dk}{(2\pi)^d} \frac{(d-2)k^\mu k^\nu+k^2\delta^{\mu\nu}}{k^2 k^2}\text{exp}\left( -2s 
k^2 \right)\nonumber\\
&=&\frac{\alpha_s}{4\pi}C_A\delta^{\mu\nu}\left[\frac{3}{\epsilon_{\rm UV}} + 3\log\left(2\mu^2 te^{\gamma_E}\right)+ \frac{5}{2}\right],
\end{eqnarray}
and
\begin{eqnarray}
\label{eq:d3}
\mathcal{M}_{(d_3)}
&=& 2(1-d)g^2C_A\tilde{\mu}^{2\epsilon}\delta^{\mu\nu}\int_0^t ds \int\frac{d^dk}{(2\pi)^d} \frac{1}{k^2}\text{exp}\left( -2s 
k^2 \right)\nonumber\\
&=&-\frac{\alpha_s}{4\pi}C_A\delta^{\mu\nu}\left[\frac{3}{\epsilon_{\rm UV}} + 3\log\left(2\mu^2 te^{\gamma_E}\right)+ 1\right],
\end{eqnarray}
from which, we see that the UV divergences in diagrams $(d_2)$ and  $(d_3)$ cancel.

Diagram $(d_4)$ gives the contribution
\begin{eqnarray}
\label{eq:d4}
\mathcal{M}_{(d_4)} 
&=& 4g^2C_A\tilde{\mu}^{2\epsilon}\int_0^t ds\int_0^s du \int\frac{d^dk}{(2\pi)^d}\frac{2(d-2)k^\mu k^\nu+\delta^{\mu\nu}k^2}{k^2 }e^{-2sk^2}\nonumber\\
&=& \frac{\alpha_s}{4\pi}C_A\delta^{\mu\nu}\left[\frac{2}{\epsilon_{\rm UV}} + 2 \log\left(2\mu^2 te^{\gamma_E}\right)  -\frac{1}{2} \right].
\end{eqnarray}

Diagram $(d_5)$ does not contribute in the small flow-time limit because it gives contribution proportional to $t$ in the small flow-time expansion.
The results of $\mathcal{M}_{(d_2)}, \mathcal{M}_{(d_3)}$ and $\mathcal{M}_{(d_4)}$ are consistent with those given in ref.~\cite{Makino:2014taa}. 

Adding up $\mathcal{M}_{(d_1)}, \mathcal{M}_{(d_2)}, \mathcal{M}_{(d_3)}$ and $\mathcal{M}_{(d_4)}$, we obtain
\begin{eqnarray}
\label{eq:dsum}
\mathcal{M}_{(d)} 
&=&\mathcal{M}_{(d_1)} + \mathcal{M}_{(d_2)} + \mathcal{M}_{(d_3)} + \mathcal{M}_{(d_4)} \nonumber\\
&=& \frac{\alpha_s}{4\pi}\delta^{\mu\nu}\left[\frac{1}{\epsilon_{\rm UV}} \beta_0 - \frac{1}{\epsilon_{\rm IR}} \left(\frac{5}{3}C_A - \frac{2}{3}n_f\right)+ 2 C_A\log\left(2\mu^2 te^{\gamma_E}\right)  +C_A\right],
\end{eqnarray}
with $\beta_0=\frac{11}{3}C_A - \frac{2}{3}n_f$.
The IR poles are removed through matching procedure and the UV poles are absorbed into the renormalization of the strong coupling in the $\overline{\rm MS}$ scheme,
\begin{eqnarray}
Z_g = 1 - \frac{\alpha_s}{4\pi}\frac{1}{2\epsilon_{\rm UV}}\beta_0 + O(\alpha_s^2),
\end{eqnarray}
which is consistent with the general argument that the flowed gluon field in $B_{\mu}$ does not need renormalization while the strong coupling in $B_{\mu}$ needs to be renormalized \cite{Luscher:2011bx}.

Finally, we have the matching coefficient for the matching from the flowed $B_{\mu}$ field to the $\overline{\rm MS}$ renormalized gluon field $gA_{\mu}$, 
\begin{eqnarray}
\label{eq:ZAresults}
\zeta_A = 1+  \frac{\alpha_s}{4\pi}C_A\left[ \log\left(2\mu^2 te^{\gamma_E}\right)  +\frac{1}{2}\right] + O(\alpha_s^2).
\end{eqnarray}
Note that the logarithmic dependence in $\zeta_A $ cancel that in $\zeta_{h_v}$ (see eq.~(\ref{eq:Zhvresult})) in adjoint representation at one-loop level. However, the constant terms in the square brackets in eq.~(\ref{eq:ZAresults}) are not canceled.
%%%%%%%%%%%%%%%%%%%%%%%%%%%%%%%%%%%%%%%
\subsection{$\zeta_{\perp\perp}$ \& $\zeta_F$}
\begin{figure}
\centering
\includegraphics[width =0.8\textwidth]{images/gluonvertexcombine2.pdf}
\caption{Vertex type Feynman diagrams that contribute to the one-loop matching for the operators $gF^{\mu\nu}_{\|\perp}h_v,\, gF^{\mu\nu}_{\perp\perp}h_v$ and $g\bar{h}_vF^{\mu\nu}_{\perp\perp}h_v$ in the gradient-flow formalism (mirror diagrams are implicit). The gluon fields at the cross circles and the filled squares are flowed at flow time $t$. For $g\bar{h}_vF^{\mu\nu}_{\perp\perp}h_v$,  the fields $\bar{h}_v$ are not drawn in diagrams $(e_1) - (e_5)$  for simplicity. The  ``heavy” auxiliary fields in $g\bar{h}_vF^{\mu\nu}_{\perp\perp}h_v$ are in fundamental representation and the ``heavy” auxiliary fields in $gF^{\mu\nu}_{\|\perp}h_v,\, gF^{\mu\nu}_{\perp\perp}h_v$ are in adjoint representation. Diagrams $(e_6)$ and $(e_7)$ only give contributions to the operators $g\bar{h}_vF^{\mu\nu}_{\perp\perp}h_v$  and $gF^{\mu\nu}_{\|\perp}h_v$, respectively.}
\label{fig:gluonvertex}
\end{figure}
The one-loop vertex type Feynman diagrams for the operators $gF^{\mu\nu}_{\|\perp}h_v,\, gF^{\mu\nu}_{\perp\perp}h_v$ and $g\bar{h}_vF^{\mu\nu}_{\perp\perp}h_v$ in the gradient-flow formalism are shown in figure \ref{fig:gluonvertex}. Choosing the external gluon momentum to be $q$ (flowing into the vertex), we have the leading-order amplitudes 
\begin{eqnarray}
\mathcal{M}_{F_{\|\perp}h_v}^{\rm LO} &=& ig\left(q_{\|}^{\mu}A^{\nu}_{\perp} -q_{_\perp}^{\nu}A^{\mu}_{\|} \right),\\
\mathcal{M}_{F_{\perp\perp }h_v}^{\rm LO} &=& ig\left(q_{\perp}^{\mu}A^{\nu}_{\perp} -q_{\perp}^{\nu}A^{\mu}_{\perp} \right),\\
\mathcal{M}_{\bar{h}_vF_{\perp\perp }h_v}^{\rm LO} &=& ig\left(q_{\perp}^{\mu}A^{\nu}_{\perp} -q_{\perp}^{\nu}A^{\mu}_{\perp} \right),
\end{eqnarray}
where we have suppressed the color indices and dropped the external ``heavy” fields but kept the external gluon fields for the convenience of matching procedure.


Since we are only interested in obtaining the matching coefficients, which are independent of the external momenta, we can set $q\cdot v=0$ for the amplitudes of the operators $gF^{\mu\nu}_{\perp\perp}h_v$ and $g\bar{h}_vF^{\mu\nu}_{\perp\perp}h_v$, but not for the amplitudes of the operator $gF^{\mu\nu}_{\|\perp}h_v$.
The computations of the one-loop vertex corrections for the operators  $gF^{\mu\nu}_{\perp\perp}h_v$ and $g\bar{h}_vF^{\mu\nu}_{\perp\perp}h_v$ are straightforward and similar. These three operators share the same diagrams $(e_1) - (e_5)$. We will only give the amplitudes of  diagrams $(e_1) - (e_5)$ for the operators $gF^{\mu\nu}_{\|\perp}h_v$ and $gF^{\mu\nu}_{\perp\perp}h_v$, supplied with the amplitude of diagram $(e_6)$ for the operator $g\bar{h}_vF^{\mu\nu}_{\perp\perp}h_v$ and the amplitude of diagram $(e_7)$ for the operator $g\bar{h}_vF^{\mu\nu}_{\|\perp}h_v$.  

The calculation of the vertex corrections for $gF^{\mu\nu}_{\|\perp}h_v$ is not so straightforward due to non-trivial mixing with the gauge non-invariant operator $iv^{\mu}A_{\perp}^{\nu}(iv\cdot D - i\delta m) h_v$.
We will deal with $gF^{\mu\nu}_{\|\perp}h_v$ in the next subsection in details. 

The amplitudes of diagrams $(e_1) - (e_6)$ are given by
\begin{eqnarray}
\label{eq:e1}
\mathcal{M}_{(e_1)} 
&=&-g^3C_A\tilde{\mu}^{2\epsilon}\int\frac{d^dk}{(2\pi)^d}\bigg\{ik^{\mu}A^{\alpha}\frac{\delta^{\nu\alpha}\left(-q-k\right)\cdot v+v^{\alpha}\left(2q-k\right)^{\nu}+v^{\nu}\left(2k-q\right)^{\alpha}}{(k\cdot v- i0)k^2(k-q)^2}\nonumber\\
&&\, \, \, \, \, \, \, \, \, \, \, \, \, \, \, \, \, \, \, \, \, \, \, \, \, \, \, \, \, \, \, \, \, \, \, \, \, \, \, \, \, \, \, \, \, \, \, \, \, \, \, \, \,  \, \, \, \, \, \, -\left(\mu\leftrightarrow\nu\right)\bigg\}e^{-t\left[k^2+(k-q)^2\right]},
\end{eqnarray}
\begin{eqnarray}
\label{eq:e2}
\mathcal{M}_{(e_2)} 
&=&g^3C_A\tilde{\mu}^{2\epsilon}\int\frac{d^dk}{(2\pi)^d}\int_0^t ds\bigg\{ik^{\mu}\frac{v^{\alpha}\left(k-2q\right)^{\nu}+2\delta^{\alpha \nu}q\cdot v-2v^{\nu}(k-q)^{\alpha}}{(k\cdot v- i0)(k-q)^2}\nonumber\\
&&\, \, \, \, \, \, \, \, \, \, \, \, \, \, \, \, \, \, \, \, \, \, \, \, \, \, \, \, \, \, \, \, \, \, \, \, \, \, \, \, \, \, \, \, \, \, \, \, \, \, \, \, \,  \, \, \, \, \, \, \, \, \, \, \, \, \, \,  \, \, \, -\left(\mu\leftrightarrow\nu\right)\bigg\}A^{\alpha}e^{-(s+t)(k-q)^2-(t-s)k^2},\\\nonumber\\
\label{eq:e3}
\mathcal{M}_{(e_3)} 
&=&g^3C_A\tilde{\mu}^{2\epsilon}\int\frac{d^dk}{(2\pi)^d}\int_0^t ds\bigg\{ik^{\mu}\frac{\delta^{\nu\alpha}\left(q+k\right)\cdot v+2v^{\nu}(-k)^{\alpha}-2v^{\alpha}q^{\nu}}{(k\cdot v- i0)k^2}\nonumber\\
&&\, \, \, \, \, \, \, \, \, \, \, \, \, \, \, \, \, \, \, \, \, \, \, \, \, \, \, \, \, \, \, \, \, \, \, \, \, \, \, \, \, \, \, \, \, \, \, \, \, \, \, \, \, \, \, \, \, \, \, \, \, \, \, \, \, \, \, \, \, \,  -\left(\mu\leftrightarrow\nu\right)\bigg\}A^{\alpha}e^{-(t+s)k^2-(t-s)(k-q)^2},\\
\label{eq:e4}
\mathcal{M}_{(e_4)} 
&=&-\frac{3}{2}g^3C_A\tilde{\mu}^{2\epsilon}\int\frac{d^dk}{(2\pi)^d } \frac{i\left(q^{\mu}A^{\nu}-q^{\nu}A^{\mu}\right)}{k^2(k-q)^2}e^{-t\left[k^2+(k-q)^2\right]},
\end{eqnarray}
\begin{eqnarray}
\label{eq:e5}
\mathcal{M}_{(e_5)} 
&=&-g^3C_A\tilde{\mu}^{2\epsilon}\int\frac{d^dk}{(2\pi)^d }\int_0^t ds\frac{i(k+3q)^{\mu}A^{\nu}-i(k+3q)^{\nu}A^{\mu}}{k^2}e^{-(t+s)k^2-(t-s)(k-q)^2},\nonumber\\\\
\mathcal{M}_{(e_6)} 
&=&g^3\left(C_F-\frac{C_A}{2}\right)\tilde{\mu}^{2\epsilon}\int\frac{d^dk}{(2\pi)^d }\int_0^t ds\frac{i\left(q^{\mu}_{\perp}A^{\nu}_{\perp}-q^{\nu}_{\perp}A^{\mu}_{\perp}\right)}{(k\cdot v-i0)^2k^2}e^{-2tk^2}.
\end{eqnarray}
Extracting the LO expressions and setting $q=0$ thereafter, we obtain the contributions from each diagram for the operator $gF^{\mu\nu}_{\perp\perp}h_v$,
\begin{eqnarray}
\label{eq:e1perpperpsimresult}
\zeta_{\perp\perp,\, (e_1)}&=&
-\frac{\alpha_sC_A}{4\pi}\times \frac{1}{2} \left[\frac{1}{\epsilon_{\rm IR}}+\log\left(2\mu^2 te^{\gamma_E}\right) +1\right],\\
\label{eq:e2perpperp}
\zeta_{\perp\perp,\, (e_2)}
&=&0,\\
\label{eq:e3perpperpresult}
\zeta_{\perp\perp,\, (e_3)}
&=&\frac{\alpha_s}{4\pi}C_A\times\frac{1}{16},\\
\label{eq:e4perpperpresult}
\zeta_{\perp\perp,\, (e_4)}&=&\frac{\alpha_sC_A}{4\pi}\times \frac{3}{2} \left[\frac{1}{\epsilon_{\rm IR}}+\log\left(2\mu^2 te^{\gamma_E}\right) +1\right],\\
\label{eq:e5perperpresult}
\zeta_{\perp\perp,\, (e_5)}
&=&- \frac{\alpha_sC_A}{4\pi}\times\frac{25}{16}.
\end{eqnarray}
Summing up all the contributions from diagrams $(e_1) - (e_5)$ and discarding the IR poles through matching procedure, we obtain the vertex matching coefficient for the operator $gF^{\mu\nu}_{\perp\perp}h_v$,
\begin{eqnarray}
\label{eq:resultgluonvertexperpperp}
\zeta_{\perp\perp} &=& 1+  \frac{\alpha_sC_A}{4\pi}\left[\log\left(2\mu^2 te^{\gamma_E}\right) -\frac{1}{2}\right]+ O(\alpha_s^2).
\end{eqnarray}
Similarly, we have the contributions from each diagram for the operator $g\bar{h}_vF^{\mu\nu}_{\perp\perp}h_v$,
\begin{eqnarray}
\zeta_{F,\, (e_1)}   &=& - \frac{\alpha_s}{4\pi}C_A\times \frac{1}{2}\left[\frac{1}{\epsilon_{\rm IR}} +\log\left(2\mu^2 te^{\gamma_E}\right) + 1\right],\\
\zeta_{F,\, (e_2)}   &=& 0,\\
\zeta_{F,\, (e_3)} &=&\frac{\alpha_s}{4\pi}C_A\times \frac{1}{16},\\
\zeta_{F,\, (e_4)}   &=&  \frac{\alpha_s}{4\pi}C_A\times \frac{3}{2}\left[\frac{1}{\epsilon_{\rm IR}} +\log\left(2\mu^2 te^{\gamma_E}\right) + 1\right],\\
\zeta_{F,\, (e_5)}   &=& - \frac{\alpha_s}{4\pi}C_A\times \frac{25}{16},\\
\zeta_{F,\, (e_6)}  &= & \frac{\alpha_s}{4\pi}\left(C_F -\frac{C_A }{2}\right)\times  2\left[\frac{1}{\epsilon_{\rm IR}} +\log\left(2\mu^2 te^{\gamma_E}\right)\right],
\end{eqnarray}  
which lead to 
\begin{eqnarray}
\label{eq:resultgluonvertexF}
\zeta_{\perp\perp}^F&= &1+ \frac{\alpha_s}{4\pi} \left[2C_F\log\left(2\mu^2 te^{\gamma_E}\right)+\frac{1}{2}C_A\right] + O(\alpha_s^2).
\end{eqnarray}
%%%%%%%%%%%%%%%%%%%%%%%%%%%%%%%%%%%%%%%%%%%%%%%%%%%%%%%
\subsection{$\zeta_{\|\perp}$}
Diagrams $(e_4), (e_5)$ have no mixing and give the same contributions for $gF^{\mu\nu}_{\perp\perp}h_v$ and $gF^{\mu\nu}_{\|\perp}h_v$, 
\begin{eqnarray}
\label{eq:e4pperpresult}
\zeta_{\|\perp,\, (e_4)}&=&\frac{\alpha_sC_A}{4\pi}\times \frac{3}{2} \left[\frac{1}{\epsilon_{\rm IR}}+\log\left(2\mu^2 te^{\gamma_E}\right) +1\right],\\
\label{eq:e5pperpresult}
\zeta_{\|\perp,\, (e_5)}
&=&- \frac{\alpha_sC_A}{4\pi}\times\frac{25}{16}.
\end{eqnarray}
\\
Complications arise from diagrams $(e_1), (e_2), (e_3)$ and $(e_7)$, whose amplitudes are given by
\begin{eqnarray}
\label{eq:e1pperp}
\mathcal{M}_{(e_1)}^{\|\perp}
&=&ig^3C_A\tilde{\mu}^{2\epsilon}\int\frac{d^dk}{(2\pi)^d}\left[\frac{(q+k)^{\mu}_{\|} A_{\perp}^{\nu}-(2q-k)^{\nu}_{\perp} A^{\mu}_{\|}}{k^2(k-q)^2} +\frac{v^{\mu}k_{\perp}^{\nu}(2k-q)_{\perp}\cdot A_{\perp} }{(k\cdot v- i0)k^2(k-q)^2} \right]\nonumber\\
&&\, \, \, \, \, \, \, \, \, \, \, \, \, \, \, \, \, \, \, \, \, \, \, \, \, \, \, \, \, \, \, \, \, \, \, \, \, \, \, \, \, \, \, \, \, \, \, \, \, \,  \, \, \, \, \, \, \, \, \, \, \, \,  \, \, \, \, \, \, \, \, \, \, \, \,  \, \, \times e^{-t\left[k^2+(k-q)^2\right]},
\end{eqnarray}
\begin{eqnarray}
\label{eq:e2pperp}
\mathcal{M}_{(e_2)}^{\|\perp} 
&=&ig_s^3C_A\tilde{\mu}^{2\epsilon}\int\frac{d^dk}{(2\pi)^d}\int_0^t ds \left[\frac{2\left(q^{\mu}_{\|}A^{\nu}_{\perp} - q_{\perp}^{\nu}A_{\|}^{\mu}\right)}{(k-q)^2} +\frac{2v^{\mu}k_{\perp}^{\nu}(k-q)\cdot A}{(k\cdot v-i0)(k-q)^2}\right]\nonumber\\
&&\, \, \, \, \, \, \, \, \, \, \, \, \, \, \, \, \, \, \, \, \, \, \, \, \, \, \, \, \, \, \, \, \, \, \, \, \, \, \, \, \, \, \, \, \, \, \, \, \, \,  \, \, \, \, \, \, \, \, \, \, \, \,  \, \, \, \, \, \, \, \, \, \, \, \,  \, \, \times e^{-(s+t)(k-q)^2-(t-s)k^2},
\end{eqnarray}
\begin{eqnarray}
\label{eq:e3pperp}
\mathcal{M}_{(e_3)}^{\|\perp}
&=&ig_s^3C_A\tilde{\mu}^{2\epsilon}\int\frac{d^dk}{(2\pi)^d}\int_0^t ds \Bigg[\frac{(k+q)^{\mu}_{\|}A^{\nu}_{\perp}-(k+q)^{\nu}_{\perp}A^{\mu}_{\|}}{k^2}\nonumber\\
&&+\frac{(A\cdot v)\left(q^{\mu}_{\|}k^{\nu}_{\perp}-q^{\nu}_{\perp}k^{\mu}_{\|}\right)+2v^{\mu}k_{\perp}^{\nu}(k\cdot A)}{(k\cdot v-i0)k^2}\Bigg]e^{-(s+t)k^2-(t-s)(k-q)^2},
\end{eqnarray}
\begin{eqnarray}
\label{eq:e7pperp}
\mathcal{M}_{(e_7)}^{\|\perp}
&=&ig^3C_A\tilde{\mu}^{2\epsilon}\int\frac{d^dk}{(2\pi)^d}\frac{-v^{\mu}A_{\perp}^{\nu} }{(k\cdot v-i0)(k-q)^2}e^{-2t(k-q)^2}.
\end{eqnarray}
These amplitudes do not give expressions that are exactly proportional to the LO expression of the operator, $gF^{\mu\nu}_{\|\perp}h_v$ even after tensor reduction for the loop integrations. They lead to mixing with the gauge non-invariant operators $iv^{\mu}A_{\perp}^{\nu}(iv\cdot D - i\delta m) h_v$, whose LO expression is proportional to $q_{\|}^{\mu}A_{\perp}^{\nu}$, and thus it is not straightforward to disentangle  the mixing from those that give contributions to the operator $gF^{\mu\nu}_{\|\perp}h_v$. Fortunately, we can extract the contributions to the operator $gF^{\mu\nu}_{\|\perp}h_v$ through the results that are proportional to $q_{\perp}^{\nu}A_{\|}^{\mu}$. To check the mixing pattern and cancellation of linear divergences that were found in refs.~\cite{Dorn:1981wa,Zhang:2018diq,Wang:2019tgg}, we will also keep the results that are proportional to $q_{\|}^{\mu}A_{\perp}^{\nu}$ and $v_{\|}^{\mu}A_{\perp}^{\nu}/\sqrt{t}$ in the small flow-time limit. 

If we set $t=0$ in $\mathcal{M}_{(e_1)}^{\|\perp}$ and $\mathcal{M}_{(e_7)}^{\|\perp}$, we can easily see that the linear divergences coming from $\mathcal{M}_{(e_7)}^{\|\perp}$ and the second term of $\mathcal{M}_{(e_1)}^{\|\perp}$ cancel with $d=3-2\epsilon$. However, in gradient  flow scheme, the cancellation of linear divergences are much more complicated. Apparently, the  linearly divergences arising from $\mathcal{M}_{(e_7)}^{\|\perp}$ and the second term of $\mathcal{M}_{(e_1)}^{\|\perp}$ do not cancel with $d=4-2\epsilon$ and $t>0$. The cancellation of linear divergences in gradient-flow scheme will be restored when diagrams $(e_2)$ and $(e_3)$ are included.

Great simplification can be achieved for our matching calculation by adopting the well-developed method used in the matching calculation in HQET. That is, we focus on the large loop momentum region $k >> q$ and expand the amplitudes up to linear order of the external momentum $q$ before performing the loop momentum integration.
This method is valid in our case because, in the small flow-time limit, the flowed operators share the same IR behavior with the corresponding un-flowed operators and contributions to the matching coefficients come from the UV region.

Expanding the amplitudes according to $k >> q$ up to linear order of $q$ and completing the loop integrations, we obtain
\begin{eqnarray}
\label{eq:e1pperpexpand}
\mathcal{M}_{(e_1)}^{\|\perp}
&=& \frac{\alpha_s}{4\pi}gC_A\left\{ -\frac{2}{3}\frac{\sqrt{2\pi}}{\sqrt{t}}v^{\mu}A_{\perp}^{\nu} - i \ell_1 \left[\left(\frac{3}{2} + \frac{4}{d}\right)q^{\mu}_{\|} A_{\perp}^{\nu}-\frac{3}{2}q^{\nu}_{\perp} A^{\mu}_{\|}\right]\right\},
\end{eqnarray}
\begin{eqnarray}
\label{eq:e2pperpexpand}
\mathcal{M}_{(e_2)}^{\|\perp} 
&=& \frac{\alpha_s}{4\pi}gC_A\left\{ -\frac{1}{6}\frac{\sqrt{2\pi}}{\sqrt{t}}v^{\mu}A_{\perp}^{\nu} + i \left[\frac{15}{8}q^{\mu}_{\|} A_{\perp}^{\nu}-\frac{1}{8}q^{\nu}_{\perp} A^{\mu}_{\|}\right]\right\},
\end{eqnarray}
\begin{eqnarray}
\label{eq:e3pperpexpand}
\mathcal{M}_{(e_3)}^{\|\perp}
&=& \frac{\alpha_s}{4\pi}gC_A\left\{ -\frac{1}{6}\frac{\sqrt{2\pi}}{\sqrt{t}}v^{\mu}A_{\perp}^{\nu} + i \left[\frac{11}{16}q^{\mu}_{\|} A_{\perp}^{\nu}-\frac{15}{16}q^{\nu}_{\perp} A^{\mu}_{\|}\right]\right\},
\end{eqnarray}
\begin{eqnarray}
\label{eq:e7pperpexpand}
\mathcal{M}_{(e_7)}^{\|\perp}
&=& \frac{\alpha_s}{4\pi}gC_A\left\{\frac{\sqrt{2\pi}}{\sqrt{t}}v^{\mu}A_{\perp}^{\nu} + \frac{i}{2} \ell_1\frac{1}{2}q^{\mu}_{\|} A_{\perp}^{\nu}\right\},
\end{eqnarray}
where 
\begin{eqnarray}
\ell_1 = \frac{1}{\epsilon_{\rm IR}}+\log\left(2\mu^2 te^{\gamma_E}\right) +1.
\end{eqnarray}
From the above results, we can clearly see the cancellation of the linear divergences. After subtracting the contributions to the operator $gF^{\mu\nu}_{\|\perp}h_v$, the leftover logarithmic dependence on $t$ also matches the UV divergence of the mixing presented in refs.~\cite{Dorn:1981wa,Zhang:2018diq,Wang:2019tgg}.

Summarizing all the contributions from diagrams $(e_1) - (e_5)$ and $(e_7)$, we obtain the matching coefficient contributed from the vertex correction of the operator $gF^{\mu\nu}_{\|\perp}h_v$,
\begin{eqnarray} 
\label{eq:resultgluonvertexpperp}
\zeta_{\|\perp} &=& 1 -\frac{1}{2}\frac{\alpha_s}{4\pi}C_A+ O(\alpha_s^2),
\end{eqnarray}
which cancels the constant term in the square brackets of $\zeta_{A}$ in eq.~(\ref{eq:ZAresults}) and thus makes the one-loop contribution to the matching coefficient of the operator $gF^{\mu\nu}_{\|\perp}h_v$ vanish \footnote{In ref.~\cite{Eller:2021qpp}, the author calculated the color-electric field correlator $G_E$ at one-loop level within the gradient-flow formalism. The matching calculation for $G_E$ in the small flow-time limit can be reduce to the matching calculation for the local current operator $g\bar{h}_vF^{\mu\nu}_{\|\perp}h_v$ with $h_v$ being in fundamental representation, which also has vanishing one-loop contribution to the matching coefficient.}.
 

%%%%%%%%%%%%%%%%%%%%%%%%%%%%%%%%%%%%%%%%%%%%%%%%%%%%%%%
\subsection{Results and applications for Wilson-line operators}
We are now ready to summarize the final results of the matching coefficients for the four local current operators: $\bar{\psi} h_v$, $gF^{\mu\nu}_{\|\perp}h_v$, $gF^{\mu\nu}_{\perp\perp}h_v$ and $g\bar{h}_vF^{\mu\nu}_{\perp\perp}h_v$. The relationships between the matching coefficients of the local current operators and the matching coefficients of the fields and vertices are provided in eqs.~(\ref{eq:operatorrelations}).
The one-loop results of $\zeta_{h_v}$, $\mathring{\zeta}_{\psi}$,  $\zeta_{\psi h_v}$, $\zeta_{A}$, $\zeta_{\|\perp}$,  $\zeta_{\perp\perp}$ and $\zeta_F$ can be found in eq.~(\ref{eq:Zhvresult}), eq.~(\ref{eq:quarkselfresult}), eq.~(\ref{eq:Zqhvresults}), eq.~(\ref{eq:ZAresults}), eq.~(\ref{eq:resultgluonvertexpperp}), eq.~(\ref{eq:resultgluonvertexperpperp}) and eq.~(\ref{eq:resultgluonvertexF}), respectively. Implementing those results, we finally obtain
\begin{eqnarray}
\label{eq:cpsihvfinal}
c_{\psi h_v}(t, \mu) &=& 1- \frac{\alpha_s}{4\pi}C_F\left[\frac{3}{2}\log\left(2\mu^2 te^{\gamma_E}\right)+\frac{\log(432)}{2}+1\right] +O(\alpha_s^2),\\
\label{eq:cpperpfinal}
c_{\|\perp} (t, \mu) &=& 1+O(\alpha_s^2),\\
\label{eq:cperpperpfinal}
c_{\perp\perp} (t, \mu) &=&1+ \frac{\alpha_s}{4\pi}C_A\times\log\left(2\mu^2 te^{\gamma_E}\right) +O(\alpha_s^2),\\
\label{eq:cpperpFfinal}
c_F(t, \mu)  &=& 1+\frac{\alpha_s}{4\pi}C_A\times\log\left(2\mu^2 te^{\gamma_E}\right) +O(\alpha_s^2).
\end{eqnarray}
We have  verified that the logarithmic dependences on $t$ in eq.~(\ref{eq:cpsihvfinal}), eq.~(\ref{eq:cpperpfinal}) and eq.~(\ref{eq:cperpperpfinal}) align with the anomalous dimensions given by refs.~\cite{Dorn:1980hs,Dorn:1981wa,Dorn:1986dt} for the current operators $\bar{\psi} h_v$, $gF^{\mu\nu}_{\|\perp}h_v$ and $gF^{\mu\nu}_{\perp\perp}h_v$, respectively (see ref.~\cite{Braun:2020ymy}  that uses similar notation for these operators for comparison). For the current operator $g\bar{h}_vF^{\mu\nu}_{\perp\perp}h_v$, its anomalous dimension  was also calculated in refs.~\cite{Eichten:1990vp,Falk:1990pz} using background field method (see ref.~\cite{Abbott:1981ke}  for a review of background field method), which is consistent with the logarithmic dependences on $t$  in eq.~(\ref{eq:cpperpFfinal}). It is noteworthy that while $gA$ does not contribute to the anomalous dimension of $g\bar{h}_vF^{\mu\nu}_{\perp\perp}h_v$ within the background field method, it does so using the method adopted in this work. The UV poles in the $\overline{\rm MS}$ scheme for diagrams $(e_1)$ and $(e_4)$ differ depending on whether the calculations are conducted with or without the background field method, owing to the distinct Feynman rules of three-gluon vertex. We have checked that these differences compensate for those arising from the anomalous dimensions of $gA$ with or without the background field method using Feynman gauge. With  correct logarithmic dependences on $t$ for the matching coefficients of  local current operators,  ``heavy” auxiliary field $h_v$,  quark field $\psi$ and $gA$, the corresponding logarithmic dependences on $t$ for the vertex matching coefficients are ensured to be correct. While the results of $c_{\perp\perp}(t, \mu)$ and $c_F(t, \mu)$ coincide at  one-loop level, it is intriguing that  the anomalous dimensions of $gF^{\mu\nu}_{\perp\perp}h_v$ and $g\bar{h}_vF^{\mu\nu}_{\perp\perp}h_v$ differ by a term proportional to $\pi^2$ at  two-loop level (see the two-loop anomalous dimensions given in refs.~\cite{Braun:2020ymy, Amoros:1997rx,Czarnecki:1997dz}). In addition, for the operator $gF^{\mu\nu}_{\|\perp}h_v$, its anomalous dimension starts to contribute at two-loop level \cite{Braun:2020ymy}, the two-loop and three-loop anomalous dimensions of the operator $\bar{\psi} h_v$ can be found in refs.~\cite{Broadhurst:1991fz,Ji:1991pr} and ref.~\cite{Chetyrkin:2003vi}, respectively.

As mentioned in the end of section \ref{subsec:smallflowexp}, the matching coefficients for the Wilson-line operators can be directly obtained from the corresponding matching coefficients of the local current operators in the small flow-time limit within the one-dimensional auxiliary-field formalism. In the small flow-time limit, expressing the matching equations for the Wilson-line operators defined in eq.~(\ref{eq:defqoperator}), eq.~(\ref{eq:defpperp}) and eq.~(\ref{eq:defperpperp}) as
\begin{eqnarray}
\mathcal{O}_{\Gamma}^{\rm R}(zv, t) &=& \mathcal{C}_{\psi}(t, \mu)e^{\delta m z}\mathcal{O}_{\Gamma}^{\overline{\rm MS}}(zv) + O(t) ,\\
\mathcal{O}_{\|\perp}^{\rm R} (zv, t) &=&  \mathcal{C}_{\|\perp}(t, \mu)e^{\delta m z}\mathcal{O}_{\|\perp}^{\overline{\rm MS}} (zv)+ O(t),\\ 
\mathcal{O}_{\|\perp}^{\rm R} (zv, t) &=& \mathcal{C}_{\perp\perp}(t, \mu)e^{\delta m z}\mathcal{O}_{\perp\perp}^{\overline{\rm MS}} (zv)+ O(t),
\end{eqnarray}
where we have omitted the Lorentz indices for simplicity, we establish relations
\begin{eqnarray}
\mathcal{C}_{\psi}(t, \mu) &=& c^2_{\psi h_v}(t, \mu),\\
\mathcal{C}_{\|\perp}(t, \mu) &=& c^2_{\|\perp} (t, \mu),\\
\mathcal{C}_{\perp\perp}(t, \mu) &=& c^2_{\perp\perp} (t, \mu).
\end{eqnarray}
These relations are independent of $z$ and external states of the matrix elements constructed using these Wilson-line operators. 

Regarding the spin-dependent potentials in terms of chromomagnetic field insertions into a Wilson loop, each such insertion will lead to a factor of the matching coefficient $c_F(t, \mu)$ for the matching from the gradient-flow scheme to the $\overline{\rm MS}$ scheme.

